\section{Identifiability, event reconstruction and modeling assumptions}
\label{sec:identifiability}

The model is expressive, and expressiveness buys ambiguity. A persistent shift
in the rate of a label can be written as a regime effect, as a slow drift of
the latent factor, or as a lagged self edge; a block of the hierarchy can be
filled in the same way by many different allocation scores. Part of that
ambiguity is statistical and shrinks as the record grows. Part of it is exact
and survives any sample size, because the parameterization contains directions
along which the observed distribution does not move at all. The two kinds have
to be handled differently. The first is a question of power and is settled by
screening a scenario before a metric is read from it; the second is a question
of what a metric is entitled to measure, and no amount of data settles it.

This section fixes the finest-scale conventions that turn the bottom of the
hierarchy into an event log, states the exact invariances as a proposition,
derives from them the quantities a parameter-recovery metric may compare, and
collects the assumptions, failure modes and safeguards under which the learned
structure and parameters carry an interpretation.

\subsection{Durations in finest-scale slots}
\label{sec:slot_durations}

Let $\Delta_1 > 0$ be the wall-clock length of one epoch at the finest scale,
called a slot. Every duration in the hierarchy is expressed as a whole number
of slots,
\begin{align}
Y_t^{(k,\ell)}
:=
\frac{\text{total wall-clock length of the episodes of label } \ell
      \text{ that start in epoch } t \text{ at scale } k}{\Delta_1}
\;\in\; \sN_0 ,
\label{eq:slot_normalized_duration}
\end{align}
so that the wall-clock duration is recovered as $\Delta_1 Y_t^{(k,\ell)}$. An
episode is booked in full to the epoch in which it starts, which is what makes
\cref{eq:slot_normalized_duration} additive over the children of a block even
when an episode runs past the end of the epoch it started in. The convention is
not cosmetic, and three consequences of it are used throughout.

First, one unit of $Y$ denotes the same amount of time at every scale, so the
duration reconciliation constraint
$\sum_{t \in \child_k(c)} Y_t^{(k,\ell)} = Y_c^{(k+1,\ell)}$
is a statement about integers. A multinomial split of an integer total
satisfies it exactly and almost surely, where a continuous duration model would
have to be corrected afterwards by a projection, as hierarchical forecast
reconciliation does \citep{wickramasuriya2019mint}.

Second, the entire model becomes discrete. Durations are then allocated by the
same Dirichlet-multinomial mechanism as counts, over the same child blocks and
with the same softmax scores, which is why the invariances of
\cref{prop:invariances} apply to the count and the duration branch in identical
form.

Third, at the finest scale a duration is an extent rather than a summary. An
episode that starts in slot $t$ and lasts $y$ slots occupies slots $t$ through
$t + y - 1$, so the reconstruction of \cref{sec:event_reconstruction} needs no
rounding step and introduces no error of its own.

Write $n_k$ for the number of finest-scale slots inside one epoch at scale $k$,
as in \cref{sec:model_scales}, so that $n_1 = 1$ and $n_K$ is the largest block
in the hierarchy. Once episodes are atomic in the sense of
\cref{assm:atomic_resolution}, an epoch at scale $k$ can start at most one
episode of a label per slot, which gives the standing bound
$X_t^{(k,\ell)} \le n_k$ on counts. Durations carry no matching per-epoch
bound. A duration is booked in full to the epoch in which its episode starts,
so $Y_t^{(k,\ell)}$ measures an extent that may run past the end of epoch $t$
and is bounded by the parent total it was allocated from,
$Y_t^{(k,\ell)} \le Y_{\parmap_k(t)}^{(k+1,\ell)}$. Iterating that bound to the
top and using the coarsest-scale ceiling of \cref{sec:atomic_resolution} gives
$Y_t^{(k,\ell)} \le n_K$ at every scale, which is a statement about the record
and not about one epoch. In particular $Y_t^{(1,\ell)}$ is the length of the
episode starting in slot $t$ and is routinely larger than one.

\subsection{Atomic resolution and the capacity cap}
\label{sec:atomic_resolution}

\begin{assumption}[Atomic resolution at the finest scale]
\label{assm:atomic_resolution}
The finest scale is chosen fine enough that a slot starts at most one episode
of any given label,
\[
X_t^{(1,\ell)} \in \{0,1\}
\qquad
\text{for all } t \in \gT^{(1)} \text{ and all } \ell \in \gL .
\]
\end{assumption}

\Cref{assm:atomic_resolution} is a statement about the choice of $\Delta_1$,
not about the mechanism that generates behavior. Under it, $X_t^{(1,\ell)}$
stops being a count and becomes an indicator that a new episode of label $\ell$
begins in slot $t$, which is what makes the finest scale reconstructible as an
event log. If the assumption fails for some label, the chosen resolution is too
coarse for that label, and the remedy is either a finer $\Delta_1$ or a
label-specific micro-allocation step below the finest scale. The assumption is
not free for the generative model either, since a Dirichlet-multinomial split
of a parent total places no bound on how much of that total lands in one child.
Two devices enforce it.

\paragraph{Capacity cap with redistribution.}
Fix a scale $k < K$, a parent epoch $c \in \gT^{(k+1)}$ and a label $\ell$, and
let $(\omega_t^{(k,\ell)})_{t \in \child_k(c)}$ be the softmax allocation
weights of \cref{eq:model_count_omega}. Every child carries the capacity $n_k$.
Write $(\tilde{X}_t^{(k,\ell)})_{t \in \child_k(c)}$ for the drawn allocation
and
\begin{align}
E_c^{(k,\ell)}
=
\sum_{t \in \child_k(c)}
\max\!\left( \tilde{X}_t^{(k,\ell)} - n_k, \; 0 \right)
\label{eq:capacity_excess}
\end{align}
for the part of it that does not fit. The draw is clipped,
\begin{align}
X_t^{(k,\ell)} = \min\!\left( \tilde{X}_t^{(k,\ell)}, \; n_k \right),
\label{eq:capacity_clip}
\end{align}
and the $E_c^{(k,\ell)}$ removed units are then reassigned one at a time, each
unit going to a child drawn from those that still have room,
\begin{align}
\Pr\!\left[\text{unit goes to } t\right]
=
\frac{\omega_t^{(k,\ell)} \,
      \mathbb{I}\{X_t^{(k,\ell)} < n_k\}}
     {\sum_{s \in \child_k(c)} \omega_s^{(k,\ell)} \,
      \mathbb{I}\{X_s^{(k,\ell)} < n_k\}} ,
\label{eq:capacity_redistribution}
\end{align}
with $X^{(k,\ell)}$ updated before the next unit is placed. The operation moves
units between siblings of one block and creates none, so the block total of
\cref{eq:model_reconciliation} holds after the repair exactly as it held
before it, and reconciliation across the whole hierarchy is untouched. The
receiver is drawn along the same allocation weights that produced the draw
being repaired, so the repair expresses no preference the model does not
already carry. It is well posed whenever the parent total fits in the block,
$X_c^{(k+1,\ell)} \le |\child_k(c)| \, n_k = n_{k+1}$, which is the count bound
of \cref{sec:slot_durations}. At the coarsest scale there is no parent total to
preserve, so the cap is applied directly to the emission draw, which truncates
the count and the duration emission at $n_K$ rather than moving units between
epochs.

\paragraph{Duration floor.}
An episode occupies at least the slot in which it starts, so before the
duration allocation distributes anything, every active child
$t \in \gA_c^{(k,\ell)}$ of \cref{eq:model_count_omega} receives a floor of
$X_t^{(k,\ell)}$ slot units, and the allocation splits only the remainder
$Y_c^{(k+1,\ell)} - \sum_{t \in \gA_c^{(k,\ell)}} X_t^{(k,\ell)}$.
The floor keeps $Y_t^{(k,\ell)} \ge X_t^{(k,\ell)}$ at every scale, which is
exactly the condition under which the next finer scale can apply the same
floor, and at the finest scale it is the condition under which every start
carries a nonempty episode. The coarsest-scale draw does not guarantee
$Y_t^{(K,\ell)} \ge X_t^{(K,\ell)}$ on its own, since the duration mean is
floored at one slot but the negative binomial draw around it is not floored at
the count. Where a parent block therefore holds fewer duration units than its
own count, the floor cannot be met, and the total is instead split across the
active children in proportion to their counts, which again preserves the block
total.

\begin{remark}
Both devices act inside one block and preserve its total. Reconciliation in
this model is a property of the generative mechanism rather than a constraint
imposed on its output, and neither the capacity cap nor the duration floor
weakens that property.
\end{remark}

\subsection{Event reconstruction as a discrete marked process}
\label{sec:event_reconstruction}

Under \cref{assm:atomic_resolution} the finest-scale arrays determine an event
log. Identify the finest epochs with the half-open wall-clock intervals
$I_t = [(t-1)\Delta_1,\, t\Delta_1)$, take the start set $\gB^{(\ell)}$ of
\cref{eq:model_episode}, and define the episode length in slots as
\begin{align}
L_t^{(\ell)} := \max\!\left( Y_t^{(1,\ell)}, \, 1 \right),
\qquad t \in \gB^{(\ell)} .
\label{eq:start_set_and_length}
\end{align}
The floor of one slot in \cref{eq:start_set_and_length} is the reconstruction
counterpart of the duration floor of \cref{sec:atomic_resolution}. A start with
zero booked duration units can arise when a parent epoch carries fewer units
than it has active children, and an episode of zero extent is not an event.

Each start induces an episode, written as a label with a start slot and an end
slot,
\begin{align}
e_t^{(\ell)}
=
\left(
\ell, \;
t, \;
f_t^{(\ell)}
\right),
\qquad
f_t^{(\ell)} = \min\!\left( t + L_t^{(\ell)} - 1, \; T_1 \right),
\qquad t \in \gB^{(\ell)},
\label{eq:episode}
\end{align}
where the end slot is inclusive and is clipped at the end of the record.
\Cref{eq:episode} is \cref{eq:model_episode} made exact at its two boundary
cases, a start that received no duration unit and an episode that would outrun
the record. The reconstructed trace of label $\ell$ is the sequence
$\gE^{(\ell)} = \langle e_t^{(\ell)} : t \in \gB^{(\ell)} \rangle$ in
increasing order of start slot, and the reconstruction map
\begin{align}
\gR :
\left\{ X_t^{(1,\ell)}, Y_t^{(1,\ell)} \right\}_{t,\ell}
\;\longmapsto\;
\left\{ \gE^{(\ell)} \right\}_{\ell \in \gL}
\label{eq:reconstruction_map}
\end{align}
is deterministic. Because its arguments are produced by exact top-down
allocation, the number of reconstructed episodes of a label under any coarser
epoch equals that epoch's count, and the reconstructed lengths sum to that
epoch's duration except where \cref{eq:start_set_and_length} lifts a zero to
one slot or \cref{eq:episode} clips at the horizon.

The traces induce the finest-scale occupancy process
\begin{align}
O_s^{(\ell)}
:=
\sum_{t \in \gB^{(\ell)}}
\mathbb{I}\!\left\{ t \le s \le f_t^{(\ell)} \right\},
\qquad s \in \gT^{(1)},
\label{eq:occupancy}
\end{align}
which counts how many episodes of label $\ell$ are active in slot $s$. Some
labels admit at most one active episode at a time, and for those the
reconstruction may return $O_s^{(\ell)} > 1$, since nothing in the allocation
prevents a long episode from running past the next start of the same label.
Collect such labels in $\gL_{\mathrm{ex}} \subseteq \gL$ and require
$O_s^{(\ell)} \le 1$ for every slot $s$ and every $\ell \in \gL_{\mathrm{ex}}$.

The requirement is enforced by a deterministic projection that truncates an
episode at the start of the next one. Write the start slots of label $\ell$ as
$b_1 < b_2 < \dots < b_{N_\ell}$, so that $N_\ell = |\gB^{(\ell)}|$. The
projected end slots are
\begin{align}
\hat{f}_{b_j}^{(\ell)}
&=
\max\!\left(
b_j, \;
\min\!\left( f_{b_j}^{(\ell)}, \; b_{j+1} - 1 \right)
\right),
\qquad j = 1,\dots,N_\ell-1,
\label{eq:exclusivity_projection}
\\
\hat{f}_{b_{N_\ell}}^{(\ell)}
&=
f_{b_{N_\ell}}^{(\ell)},
\label{eq:exclusivity_projection_last}
\end{align}
with the outer maximum keeping every projected episode at least one slot long
and the horizon clip inherited from \cref{eq:episode}.

\begin{remark}
\label{rem:projection_is_a_safeguard}
The projection changes no start slot and removes no episode, so the count of
episodes per label is preserved; it only shortens durations. It is the third
place, alongside the one-slot floor and the horizon clip, where the
reconstruction can depart from exact duration reconciliation, because the slots
it removes are not returned to a sibling. It is therefore a
safeguard rather than a mechanism, and the burden of avoiding overlap belongs
to the model, through the lagged structure, the context features and the
label-specific constraints. Evaluation reflects both of the things the
projection touches, scoring start slots by detection $F_1$ against the true
start set and extents by the intersection over union of the occupied slot sets
$\{ s \in \gT^{(1)} : O_s^{(\ell)} > 0 \}$ of \cref{eq:occupancy}.
\end{remark}

\subsection{Fundamental invariances}
\label{sec:invariances}

The invariances below are exact. Along each of them the joint distribution of
the observed counts and durations given the context is literally unchanged, so
no estimator, no sample size and no regularizer can distinguish the two
parameter values, and any evaluation that appears to distinguish them is
measuring its own coordinate system.

\begin{proposition}[Exact invariances of the parameterization]
\label{prop:invariances}
Fix a scale $k$ and let $\Theta'$ be obtained from $\Theta$ by any one of the
following transformations. Then $\Theta$ and $\Theta'$ induce the same
distribution of $(X,Y)$ given the context, that is, they are observationally
equivalent.
\begin{enumerate}
\item \textbf{Regime relabeling.} For a permutation $\sigma$ of
$\{1,\dots,M_k\}$, replace $S_t^{(k)}$ by $\sigma(S_t^{(k)})$ and permute every
block indexed by a scale-$k$ regime consistently, so that $a_m^{(k)} \mapsto
a_{\sigma^{-1}(m)}^{(k)}$, $A_{m,m'}^{(k)} \mapsto
A_{\sigma^{-1}(m),\sigma^{-1}(m')}^{(k)}$, and likewise for the context weights
$\vb_m^{(k)}$, the initial logits $\iota_m^{(k)}$, the factor blocks
$(\mF_m^{(k)}, \mG_m^{(k)}, \mSigma_m^{(k)})$ and the regime effects
$\vbeta_{S,\ell}^{(k)}$, $\valpha_{S,\ell}^{(k)}$,
$\boldsymbol{\delta}_{S,\ell}^{(k)}$. The cross-scale matrix $C^{(k)}$ within
the transition map of \cref{eq:model_regime_init} is indexed by a scale-$k$ regime in its first slot
and by a scale-$(k+1)$ regime in its second, so $\sigma$ permutes the first
index of $C^{(k)}$ and the second index of $C^{(k-1)}$, and the invariance
holds one scale at a time and, applied at every scale, jointly.

\item \textbf{Factor rotation and sign.} For any invertible
$\mR \in \sR^{r_k \times r_k}$, replace $H_t^{(k)}$ by $\mR H_t^{(k)}$, the
factor dynamics by $(\mR \mF_m^{(k)} \mR^{-1}, \, \mR \mG_m^{(k)}, \,
\mR \mSigma_m^{(k)} \mR^\top)$ for every $m$, the factor initialization by
$(\mR \vmu_0^{(k)}, \mR \mSigma_0^{(k)} \mR^\top)$, and every loading vector by
$\vbeta_{H,\ell}^{(k)} \mapsto \mR^{-\top} \vbeta_{H,\ell}^{(k)}$, and
similarly for $\valpha_{H,\ell}^{(k)}$ and $\boldsymbol{\delta}_{H,\ell}^{(k)}$.
Writing $\mB_H^{(k)} \in \sR^{r_k \times L}$ for the matrix whose $\ell$-th
column is $\vbeta_{H,\ell}^{(k)}$, only the row space of $\mB_H^{(k)}$ is
pinned down, and in particular a sign flip of a factor coordinate is invisible.

\item \textbf{Block-constant shifts of an allocation score, $k < K$.} For
constants $(v_\ell)_{\ell \in \gL}$, any one of the shifts
\begin{align*}
\beta_{0,\ell}^{(k)} \mapsto \beta_{0,\ell}^{(k)} + v_\ell,
\qquad
&\beta_{P,\ell}^{(k)} \mapsto \beta_{P,\ell}^{(k)} + v_\ell,
\qquad
\beta_{D,\ell}^{(k)} \mapsto \beta_{D,\ell}^{(k)} + v_\ell,
\\
&\beta_{S,m,\ell}^{(k)} \mapsto \beta_{S,m,\ell}^{(k)} + v_\ell
\;\;\text{for every } m,
\end{align*}
leaves every count allocation probability of \cref{eq:model_count_omega}
unchanged, and the corresponding shifts of $\alpha_{0,\ell}^{(k)}$,
$\alpha_{P,\ell}^{(k)}$ and $\alpha_{S,m,\ell}^{(k)}$ leave the duration
allocation produced by \cref{eq:model_count_omega} unchanged.

\item \textbf{Intercept trade at the coarsest scale.} For constants
$(v_\ell)_{\ell \in \gL}$, the simultaneous shift
\[
\beta_{0,\ell}^{(K)} \mapsto \beta_{0,\ell}^{(K)} + v_\ell,
\qquad
\beta_{S,m,\ell}^{(K)} \mapsto \beta_{S,m,\ell}^{(K)} - v_\ell
\;\;\text{for every } m,
\]
leaves the count mean unchanged, and the same trade in
$(\delta_{0,\ell}^{(K)}, \boldsymbol{\delta}_{S,\ell}^{(K)})$ and in
$(\alpha_{0,\ell}^{(K)}, \valpha_{S,\ell}^{(K)})$ leaves the zero-inflation
probability and the duration mean unchanged. Only the sums
$\beta_{0,\ell}^{(K)} + \beta_{S,m,\ell}^{(K)}$ are pinned down.

\item \textbf{Destination-free shifts of the regime logits.} For any constants
$v_0$, $(v_{m'})_{m'}$, $(v_n)_n$ and any vector $\vv \in \sR^{d_k}$, the
shifts
$a_m^{(k)} \mapsto a_m^{(k)} + v_0$,
$A_{m,m'}^{(k)} \mapsto A_{m,m'}^{(k)} + v_{m'}$,
$C_{m,n}^{(k)} \mapsto C_{m,n}^{(k)} + v_n$
and
$\vb_m^{(k)} \mapsto \vb_m^{(k)} + \vv$,
applied for every destination $m$, leave every regime transition probability
unchanged. Only the row-stochastic transition table at a given context is
pinned down.
\end{enumerate}
\end{proposition}

\begin{proof}
Item 1 is a relabeling of a summation index. The latent regime is summed out of
the observed likelihood, and permuting the label of the summand together with
every block it indexes leaves each term of the sum and therefore the sum
itself.

Item 2 holds because the factor enters every emission and every allocation
score only through the inner products $(\vbeta_{H,\ell}^{(k)})^\top H_t^{(k)}$
and their $\valpha$ and $\boldsymbol{\delta}$ analogues, and
$(\mR^{-\top}\vbeta_{H,\ell}^{(k)})^\top (\mR H_t^{(k)})
= (\vbeta_{H,\ell}^{(k)})^\top H_t^{(k)}$. The transformed dynamics are those
of the process $\mR H_t^{(k)}$, and a linear map of a Gaussian is Gaussian with
the stated covariance, so the joint law of the inner products is unchanged.
Rows of $\mR^{-\top}\mB_H^{(k)}$ are linear combinations of the rows of
$\mB_H^{(k)}$ and $\mR^{-\top}$ is invertible, so the row space is the
invariant object.

Item 3 holds because the allocation probability of a child is a softmax over
the children of one block, \cref{eq:model_count_omega}, and a softmax is
invariant to adding to every score of the block a quantity that does not depend
on the child. The intercept is such a quantity trivially. The parent terms
$v_\ell \log(1 + X_c^{(k+1,\ell)})$ and $v_\ell \log(1 + Y_c^{(k+1,\ell)})$
depend on the parent epoch $c$ but on no child of it, so they are constant
inside the block. A per-label constant added to the whole regime column
contributes $v_\ell$ for every child, whatever regime that child occupies,
because exactly one regime is active per epoch. The same argument applies
verbatim to the duration score over the active children $\gA_c^{(k,\ell)}$,
with $\alpha_{P,\ell}^{(k)}$ in place of the two parent terms. No other feature
of \cref{eq:model_count_score} is block-constant, since the
factor, the context, the depletion feature $W_t^{(k,\ell)}$, the own-epoch
count and the lagged values all differ between children of the same block.

Item 4 holds because the coarsest-scale linear predictors see the regime only
through $\beta_{0,\ell}^{(K)} + \beta_{S,m,\ell}^{(K)}$ evaluated at the active
regime $m = S_t^{(K)}$, and the trade leaves that sum, hence
$\mu_t^{(K,\ell)}$, hence the emission law; the argument for
$\varpi_t^{(K,\ell)}$ and $\nu_t^{(K,\ell)}$ is the same predictor written with
$\boldsymbol{\delta}$ and $\valpha$.

Item 5 holds because the transition probability of
\cref{eq:model_regime_init} is a softmax over the
destination regime $m$, and each listed shift adds to the logit vector a
quantity that is constant in $m$ for every fixed conditioning value of the
previous regime, the parent regime and the context. The shift of
$\vb_m^{(k)}$ is likewise constant in $m$ in the initial logits of
\cref{eq:model_regime_init}, so the initial distribution is unchanged as well.
\end{proof}

\begin{remark}
\label{rem:not_invariant}
The complement is as important. A lagged coefficient
$\beta_{\ell',d}^{(k,\ell)}$ is identified, because its regressor
$\log(1 + X_{t-d}^{(k,\ell')})$ varies across the children of a block; so is
the depletion coefficient $\beta_{\mathrm{dep},\ell}^{(k)}$, whose regressor
$\log(1 + W_t^{(k,\ell)})$ is built from what earlier children of the same
block have already taken and therefore varies inside the block by
construction; so is the count coupling $\alpha_{X,\ell}^{(k)}$, and so are the
dispersions $\phi_\ell^{(K)}$, $s_\ell^{(K)}$ and the concentrations
$\kappa_k$, $\tau_k$. The contrast of the regime effect across regimes is
identified even where its per-label level is not. These are the directions in
which a learner can be right or wrong.
\end{remark}

\subsection{Identified quantities and what a recovery metric may compare}
\label{sec:identified_quantities}

\Cref{prop:invariances} has a direct operational consequence. A parameter is
not a measurable quantity; a function of parameters that is constant on the
orbits of the invariances is. Reading off the proposition, the identified
objects per scale are the following.

\begin{itemize}
\item The row-stochastic regime transition table at a reference context, which
by item 5 is what the transition logits determine.
\item The row space of the factor loading block, by item 2.
\item At the coarsest scale, the array of sums
$\beta_{0,\ell}^{(K)} + \beta_{S,m,\ell}^{(K)}$ over regimes $m$ and labels
$\ell$, by item 4, together with the count dispersion $\phi_\ell^{(K)}$ and the
duration dispersion $s_\ell^{(K)}$.
\item At a finer scale, the regime effect after its per-label mean over regimes
is removed, by item 3. Equivalently, and more directly interpretable, the
shares that a block would give its children if those children sat in different
regimes, which is the softmax of the centered regime effect down the regime
axis.
\item The lagged coefficients on the true support, and the total absolute
lagged mass a learner places off that support.
\end{itemize}

The parameter-recovery metric is built from this list and reads nothing outside
it. It first aligns the regimes of the learned model with those of the truth by
the permutation that minimizes the squared distance between the identified
regime levels, so that a model differing only by a relabeling is aligned before
anything is scored. It then reports, per scale, the mean total variation
distance between the transition tables, the mean absolute error of the
identified emission blocks and that error divided by the spread of the true
block, one minus the mean best absolute cosine between true and learned loading
directions, the mean absolute difference of the implied allocation shares, and
the two lagged-coefficient summaries. The emission blocks are the identified
regime level at every scale and, at the coarsest scale, the count dispersion
alongside it. The absolute value in the cosine is the sign invariance of item 2
made explicit. The single headline number averages the scale-free entries
across scales.

\begin{remark}
The loading comparison uses the mean best absolute cosine rather than a
subspace distance. The two agree when $r_k = 1$ and when the transformation
relating the two loading blocks is a signed permutation; under a general
rotation the cosine statistic is a conservative summary of the row-space
agreement of item 2 rather than an exact one.
\end{remark}

The reason for this care is negative rather than positive. A metric that
compares raw parameter arrays entry by entry charges a learner for every step
it takes along a direction of \cref{prop:invariances}, which is a direction the
data cannot see. Two models that induce the same distribution of the observed
record then receive different scores, and the ranking they induce is partly a
ranking of coordinate conventions. This is the same class of failure that scale
artifacts introduce into structure benchmarks, where a criterion picks up a
property of the simulation rather than the property under study
\citep{reisach2021varsortability,reisach2023r2sortability}, and the same reason
structural comparison is done modulo an equivalence class rather than on raw
adjacency \citep{peters2015sid,runge2018pitfalls}.

The construction has an empirical anchor. In the parameter-recovery experiment
an oracle method returns the generating parameters without fitting anything,
and its headline error is zero to machine precision on all $27$ scenario and
replicate rows of \texttt{results/exp7/exp7\_results.csv}, the largest value
being $2.8 \times 10^{-17}$. The floor is exact by construction, which is what
makes the errors reported by the fitted methods above it interpretable as
error rather than as bookkeeping.

\subsection{Separation of explanatory roles}
\label{sec:role_separation}

The three mechanisms that carry temporal and cross-label dependence are the
latent regime, the latent factor and the lagged edges. They are not
interchangeable, and the model needs all three because each covers a pattern
the other two either cannot express or would express by distorting something
else.

The regime $S_t^{(k)}$ carries persistent drift. It is a discrete state with a
sticky transition table, so it changes rarely and, while it holds, shifts the
level of counts and durations for many labels at once. A record in which the
base rate of a label is high for a stretch of epochs and low afterwards is what
a regime is for.

The factor $H_t^{(k)}$ carries same-epoch co-movement. It is continuous,
low-dimensional and shared across labels within one epoch, so it explains why
two labels move together at the same time without asserting that either causes
the other. This is the mechanism that makes the model a dynamic Bayesian
network in the strict sense, because it lets all contemporaneous association
among observed variables be mediated by a latent quantity rather than by a
same-epoch directed edge.

The lagged parent sets $\gP_{k,\ell}$ carry directed dependence. An edge
$(\ell',d) \to \ell$ at scale $k$ states that the value of label $\ell'$ at lag
$d$ shifts the conditional law of label $\ell$ now, given everything else
admissible.

Dropping any one of the three does not simplify the model; it redirects the
variation into the remaining two. Without a regime, a persistent level shift is
absorbed by lagged self edges, since a persistent state is autocorrelated and a
lagged self edge is the cheapest way for a criterion that conditions on nothing
to express autocorrelation. This is the over-selection mechanism that the
results section quantifies. Without a factor, same-epoch co-movement has no
admissible carrier and is expressed as directed dependence at whatever lag the
sampling makes it look like. Without lagged edges, genuine directed dependence
is pushed onto a latent state, which then splits into more regimes than the
phenomenon has.

\begin{assumption}[Role separation]
\label{assm:role_separation}
Within a scale, parent-scale summaries explain the transfer of a coarse budget
to finer epochs, the lagged parent sets explain within-scale directed
dependence, the regime explains persistent drift in local dynamics, and the
factor explains same-epoch co-movement not attributable to lagged structure.
\end{assumption}

\Cref{assm:role_separation} is partly enforced by construction. Prohibiting
same-epoch directed edges among observed variables prevents the graph from
absorbing contemporaneous association; exact reconciliation prevents finer
scales from redefining coarser totals. The remainder is a modeling commitment,
and the ablation that removes each component in turn is what tests it.

\subsection{Modeling assumptions}
\label{sec:assumptions}

The following assumptions are what the interpretation of a fitted model rests
on. They are stated so that a reader can check them against a record before
reading a result off it.

\begin{assumption}[Correct multiscale aggregation]
\label{assm:aggregation}
The observed arrays at all scales are deterministic aggregations of one
finest-scale event log, so the reconciliation constraints hold in the data and
not only in the model.
\end{assumption}

\begin{assumption}[Consistent label semantics across scales]
\label{assm:label_semantics}
A label denotes the same phenomenon at every scale. If aggregating over a
coarser window changes what a label means, the hierarchy relates quantities
that are not comparable and the top-down construction has no referent.
\end{assumption}

\begin{assumption}[Admissible lag structure]
\label{assm:lag_structure}
All directed dependencies among observed variables are lagged and within a
scale. No directed edge joins two observed variables in the same epoch, and all
contemporaneous association is mediated by $H_t^{(k)}$ and $S_t^{(k)}$.
\end{assumption}

\begin{assumption}[Finite candidate parents and sparse truth]
\label{assm:candidates}
For each scale $k$ and label $\ell$ the candidate set
$\gC_{k,\ell} \subseteq \gL \times \{1,\dots,p_k\}$ is finite and contains the
true parent set, and the true parent set has bounded cardinality.
\end{assumption}

\begin{assumption}[Distinct regime signatures]
\label{assm:distinct_regimes}
At each scale, distinct regimes induce distinct conditional behavior in at
least one of the transition table, the factor dynamics, or the emissions of at
least one label. No two regimes are observationally identical beyond the
relabeling of item 1 of \cref{prop:invariances}.
\end{assumption}

\begin{assumption}[Factor normalization]
\label{assm:factor_normalization}
The innovation covariance is fixed at a reference regime, which cuts the
freedom of item 2 of \cref{prop:invariances} from the general linear group to
the orthogonal group. No location normalization is imposed, since
the continuous component of \cref{eq:model_regime_init} carries no intercept and the mean of the
factor is therefore already determined by $\vmu_0^{(k)}$, $\mF^{(k)}$,
$\mG^{(k)}$ and the context rather than free to be set. Recovery is claimed for
the loading row space and not for individual loadings.
\end{assumption}

\begin{assumption}[Allocation feature non-redundancy]
\label{assm:nonredundancy}
Within a child block, the features entering an allocation score are not
perfectly collinear. Beyond the block-constant shifts of item 3 of
\cref{prop:invariances}, distinct allocation parameters then induce distinct
allocation probabilities.
\end{assumption}

\begin{assumption}[Stationarity within a regime]
\label{assm:within_regime_stationarity}
Conditional on the regime, the context and the parent-scale summaries, local
dynamics are approximately stable. The model is non-stationary across regimes,
not within one.
\end{assumption}

\begin{assumption}[Regularity]
\label{assm:regularity}
Parameters lie in the interior of the parameter space and every conditional
distribution used is well defined and non-degenerate over the support of the
data.
\end{assumption}

Two remarks bound what these assumptions buy. First, exact global
identifiability of the graph is not available and is not claimed. Under
\cref{assm:candidates} the target is restricted identifiability inside the
admissible class, and an edge is recoverable when its inclusion changes the
conditional law of $(X_t^{(k,\ell)}, Y_t^{(k,\ell)})$ given all other
admissible predictors by a non-negligible amount over the support of the data.
Two nearly collinear candidate predictors are not separable by any criterion at
finite sample size, so a recovered graph is a statement about statistically
distinguishable dependencies relative to the candidate set and the penalty
$\lambda_k$, not about a unique mechanism. Second,
\cref{assm:lag_structure} makes the factor the carrier of contemporaneous
association, but a low-dimensional factor does not adjust for every unobserved
influence. Learned lagged edges are predictive temporal relations unless
stronger assumptions are added, which is the standing caveat for structure
learned from observational time series with possible latent confounding
\citep{gerhardus2020lpcmci}. Regime label switching and the resulting need to
align before comparing are shared with the non-stationary dynamic Bayesian
network literature more broadly
\citep{robinson2010nsdbn,grzegorczyk2011nhdbncontinuous}.

\subsection{Failure modes and safeguards}
\label{sec:failure_modes}

Each assumption above fails in a characteristic way, and each failure has a
signature that can be checked.

A finest scale that is too coarse violates \cref{assm:atomic_resolution}. Its
signature is finest-scale counts above one before the capacity cap acts, and
its consequence is that reconstruction stops being exact, since several
episodes of a label share a slot and their boundaries are no longer recoverable
from $(X^{(1)}, Y^{(1)})$.

Redundant candidate predictors violate \cref{assm:nonredundancy} and
\cref{assm:candidates}. Their signature is a structure score that is nearly
flat over a set of competing edges, and their consequence is that recovered
parent sets vary across replicates while predictive performance does not.

An over-flexible factor absorbs genuine lagged dependence. Raising $r_k$ gives
the factor enough freedom to reproduce, as same-epoch co-movement, a pattern
that a lagged edge would explain. The signature is a fall in the recall of the
true edges with little movement in the held-out likelihood.

Too many regimes split one behavior across several states. The signature is a
transition table with several rows that are close to each other and regimes
whose identified emission levels differ by little, which also makes the regime
alignment of \cref{sec:identified_quantities} unstable.

Misspecified allocation features confound top-down transfer with within-scale
dynamics. If the competition among the children of a block for a fixed parent
total is not represented by a feature that varies inside the block, the only
instrument the structure search has for it is a lagged self edge, and
within-block competition is then reported as structure.

The safeguards follow from the list. Keep the regime cardinalities $M_k$ and
factor dimensions $r_k$ small, and prefer the smallest values that a held-out
criterion supports. Bound the candidate lag sets and the in-degree. Penalize
graph and parameter complexity together. Screen a scenario for identifiability
before reading a metric from it, by checking that the likelihood gain of a true
edge exceeds both the gain of the best non-edge and the penalty an edge must
beat, and by comparing latent recovery against the oracle attainable with the
true parameters; \cref{sec:benchmark_identifiability} gives the screen used
here and the reading it returns. Ablate the components against each other, so
that parent-scale summaries, factor and lagged edges are each shown to carry
predictive value the others do not.

\begin{remark}
The model is a structured, partially identifiable multiscale latent-variable
model. Its parameters and its graph are meaningful relative to the admissible
class, the chosen resolutions, and the normalization and sparsity assumptions
above, and any claim read from a fitted model inherits those qualifications.
\end{remark}
